\clearpage
\section{Iterative Proportional Representation with Guaranteed Winner}

\subsection{Procedure Description}

Proportional representation tends to force coalitions because no party receives an absolute
majority. The more large parties there are, the lower the probability that
voters clearly prefer one. The iterative electoral procedure proposed below ensures
that the compulsion to form grudging coalitions with rotten compromises is eliminated. It
guarantees that a governing majority always emerges in the end. However, it also gives
voters who chose \enquote{losing parties} in the first round the opportunity in further rounds
to help determine, in the sense of the lesser evil, who will ultimately govern. This
distinguishes it from other majority-forming electoral procedures.

\begin{enumerate}
  \item In the first voting step, a normal proportional election is conducted in which the parties
    are ordered according to their percentage of votes. If a party receives an
    absolute majority or if a coalition is formed that has an absolute majority,
    the procedure is terminated and the party or coalition with the absolute majority has won the election.
    The \textbf{absolute majority} in this sense is a predefined
    percentage that is slightly above 50\,\%. It will be set so that the
    governing party/coalition gets a few seats more than half, so that it
    still has the majority mathematically even when someone is sick.

  \item If no winner is determined, two cases must be distinguished:

    \begin{enumerate}
      \item Only two parties were available for selection. This case is again divided into two
        cases:

        \begin{enumerate}[label=\roman*.]
          \item It was the first round with this composition. In this case,
            step~1 is repeated once more.

          \item It is already the second inconclusive round with only two parties. In this
            case, a random procedure decides, for example minimal percentage
            differences in the result, or drawing lots.
        \end{enumerate}

      \item More than two parties were available for selection. In this case, the percentage
        vote shares are summed up in order, starting with the largest, until a
        two-thirds majority is reached or just exceeded. This means that at least
        two-thirds of the voters want at least one of these parties to govern. Now
        step~1 is executed again, but only with the parties determined in this way. All voters are
        now allowed to vote again on these parties.
    \end{enumerate}
\end{enumerate}

The procedure is repeated until there is a winner. The procedure always terminates
because the number of parties running for election is reduced in each iteration step,
until possibly only two parties remain. The idea behind the procedure is that
people usually also make a decision when they are forced to choose between two
alternatives --- even if it's only the lesser of two evils. With many
alternatives, however, one person chooses this, another that. Coalitions, when they
come from the heart, are preferable to multiple rounds of voting, as they come closest to the will of the voters. However,
there is no longer the compulsion to form coalitions at any cost just so that someone can govern.
The parties can consider how their chances would stand in another round of voting
and possibly seek a coalition partner or not. It may well happen that
two parties are content-wise close, but are simply two different teams competing
for the favor of the voters. In this case, a coalition would also be advisable from the voters' perspective.

If in the first round a party or coalition was already able to achieve a governing majority ($\geq$
absolute majority), it receives a corresponding number of seats allocated in parliament.
Otherwise, the winner from multiple rounds of voting receives a number of parliamentary seats
corresponding to the absolute majority. The remaining seats are distributed among the
other parties according to the original ratio in the first round. This
ensures that there is always a governing majority, but also that smaller parties in the
opposition are represented, as long as their vote result is sufficient for at least one representative.

\bigskip

\runinheadul{Example:}

\medskip\noindent
\begin{tabular}{@{}>{\centering\arraybackslash}b{0.48\textwidth}@{\hspace{0.04\textwidth}}>{\centering\arraybackslash}b{0.48\textwidth}@{}}

  % --- 1st Round ---
  1st Round

  \medskip

  \begin{tikzpicture}[x=2.5cm, y=0.08cm]
    \fill[diagrammHG] (0,0) rectangle (1,100);
    \fill[parteiA] (0,0)     rectangle (1,30);
    \fill[parteiB] (0,30)    rectangle (1,52.5);
    \fill[parteiC] (0,52.5)  rectangle (1,72.5);
    \fill[parteiD] (0,72.5)  rectangle (1,82.5);
    \fill[parteiE] (0,82.5)  rectangle (1,90.75);
    \fill[parteiF] (0,90.75) rectangle (1,96.75);
    \fill[parteiS] (0,96.75) rectangle (1,100);
    \draw (0,0) rectangle (1,100);
    % Percentage values inside
    \node at (0.5,15)    {\small 30\,\%};
    \node at (0.5,41.25) {\small 22{,}5\,\%};
    \node at (0.5,62.5)  {\small 20\,\%};
    \node at (0.5,77.5)  {\small 10\,\%};
    \node at (0.5,86.6)  {\small 8{,}25\,\%};
    \node at (0.5,93.75) {\small 6\,\%};
    % Party abbreviations on the right
    \node[right] at (1,15)    {A};
    \node[right] at (1,41.25) {B};
    \node[right] at (1,62.5)  {C};
    \node[right] at (1,77.5)  {D};
    \node[right] at (1,86.6)  {E};
    \node[right] at (1,93.75) {F};
    \node[right] at (1,98.4)  {S};
    % Dimension arrow on the left (y=0 to y=66)
    \draw (0,0)  -- (-0.25,0);
    \draw (0,66) -- (-0.25,66);
    \draw[<->] (-0.25,0) -- (-0.25,66);
    \node[left] at (-0.3,33) {\small 66\,\%};
  \end{tikzpicture}

  &

  % --- 2nd Round ---
  2nd Round

  \medskip

  \begin{tikzpicture}[x=2.5cm, y=0.08cm]
    \fill[diagrammHG] (0,0) rectangle (1,100);
    \fill[parteiA] (0,0)  rectangle (1,45);
    \fill[parteiC] (0,45) rectangle (1,75);
    \fill[parteiB] (0,75) rectangle (1,100);
    \draw (0,0) rectangle (1,100);
    % Percentage values inside
    \node at (0.5,22.5) {\small 45\,\%};
    \node at (0.5,60)   {\small 30\,\%};
    \node at (0.5,87.5) {\small 25\,\%};
    % Party abbreviations on the right
    \node[right] at (1,22.5) {A};
    \node[right] at (1,60)   {C};
    \node[right] at (1,87.5) {B};
    % Dimension arrow on the left (y=0 to y=66)
    \draw (0,0)  -- (-0.25,0);
    \draw (0,66) -- (-0.25,66);
    \draw[<->] (-0.25,0) -- (-0.25,66);
    \node[left] at (-0.3,33) {\small 66\,\%};
  \end{tikzpicture}

\end{tabular}

\bigskip

\noindent
\begin{tabular}{@{}>{\centering\arraybackslash}b{0.48\textwidth}@{\hspace{0.04\textwidth}}>{\centering\arraybackslash}b{0.48\textwidth}@{}}

  % --- 3rd Round ---
  3rd Round

  \medskip

  \begin{tikzpicture}[x=2.5cm, y=0.08cm]
    \fill[diagrammHG] (0,0) rectangle (1,100);
    \fill[parteiA] (0,0)    rectangle (1,47.5);
    \fill[parteiC] (0,47.5) rectangle (1,100);
    \draw (0,0) rectangle (1,100);
    % Percentage values inside
    \node at (0.5,23.75) {\small 47{,}5\,\%};
    \node at (0.5,73.75) {\small 52{,}5\,\%};
    % Party abbreviations on the right
    \node[right] at (1,23.75) {A};
    \node[right] at (1,73.75) {C};
    % Dimension arrow on the left (y=0 to y=66)
    \draw (0,0)  -- (-0.25,0);
    \draw (0,66) -- (-0.25,66);
    \draw[<->] (-0.25,0) -- (-0.25,66);
    \node[left] at (-0.3,33) {\small 66\,\%};
  \end{tikzpicture}

  &

  % --- Result Distribution ---
  Result Distribution

  \medskip

  \begin{tikzpicture}[x=2.5cm, y=0.08cm]
    \fill[diagrammHG] (0,0) rectangle (1,100);
    \fill[parteiC] (0,0)    rectangle (1,51);
    \fill[parteiA] (0,51)   rectangle (1,69.3);
    \fill[parteiB] (0,69.3) rectangle (1,83.1);
    \fill[parteiD] (0,83.1) rectangle (1,89.2);
    \fill[parteiE] (0,89.2) rectangle (1,94.3);
    \fill[parteiF] (0,94.3) rectangle (1,98);
    \draw (0,0) rectangle (1,100);
    % Percentage values inside
    \node at (0.5,25.5)  {\small 51\,\%};
    \node at (0.5,60.15) {\small 18{,}3\,\%};
    \node at (0.5,76.2)  {\small 13{,}8\,\%};
    \node at (0.5,86.15) {\small 6{,}1\,\%};
    \node at (0.5,91.75) {\small 5{,}1\,\%};
    \node at (0.5,96.15) {\small 3{,}7\,\%};
    % Party abbreviations on the right
    \node[right] at (1,25.5)  {C};
    \node[right] at (1,60.15) {A};
    \node[right] at (1,76.2)  {B};
    \node[right] at (1,86.15) {D};
    \node[right] at (1,91.75) {E};
    \node[right] at (1,96.15) {F};
    % Phantom node for alignment (same width as dimension arrow of other diagrams)
    \node[opacity=0] at (-0.3,33) {\small 66\,\%};
  \end{tikzpicture}

\end{tabular}

\bigskip

In the above example, party~C wins the election in the 3rd round, even though A was the
strongest force in the first round. If A and B had formed a coalition after the first round,
the election would have been over and A and B would have governed together. But since they had incompatible
positions or persons, they decided to proceed to the next round.
Only as many parties participate in the second round as are just necessary to exceed two-thirds
of the votes. In this case, that was A, B, and C. In the second round, A
was able to expand its lead and now reaches 45\,\%. Interestingly, C and B have
switched places, as apparently more of the voters of the parties eliminated in the first round
were more inclined toward party~C than party~B.
Since still no party has the absolute majority and coalitions are ruled out,
the 3rd round comes down to a runoff between A and C. Interestingly, A loses its
previous leading role in the runoff, as most former B-voters vote for C. C has
thus won the election and is awarded the absolute majority. Since this is only an
awarded absolute majority due to runoffs, the exact percentage
result of the runoffs is irrelevant: The winner only receives the awarded number of seats. (In
this example 51\,\%) This rule is important because the runoffs could also end with
a much larger difference. For example, A could also receive 70\,\% of the votes.
All other parties would then be absolutely underrepresented, which would not correspond to the actual
will of the people. If in the first round a party or coalition were to unite 70\,\% of the votes,
that would be a different statement. Then 70\,\% of the population really want
this party/parties to govern, despite the multitude of other options.

Therefore, absolute majorities obtained in the first round may remain unchanged.

\subsection{Rejected Alternative}

Instead of awarding the absolute majority of votes, one could simply grant the
winner party of the iterative procedure the right to legislate. Why vote at all
anymore? They win every vote anyway since they have the absolute majority.
However, this leads to the question of why we need parliaments at all. If the factions always
vote exactly as their faction leader says anyway, we would indeed no longer need
parliaments. Then it would be sufficient to form a faction leaders' council, where the
faction leaders would vote with a certain voting weight. If the
faction leader of the governing party were also the head of government,
autocracy would be complete.
But that's not how it is. Parliamentarians can often vote by secret ballot and the
faction leaders must convince their faction members and cannot
command them. For some questions, one cannot even rely on one's own party members,
but must seek allies in other factions. This is exactly as desired,
so that the power of individuals does not become too great. Important questions should not be
decided by individuals, but by a group of people. However, one assumes that people
with the same basic convictions, as can be assumed from members of a party,
will usually come to an agreement. However, some questions are also so
controversially discussed within a party that agreement is not always possible. The larger and more heterogeneous a
faction is, the more difficult an agreement will generally be. In view of
limiting power, it therefore makes sense to build in a compulsion to reach agreement among more people
as compensation for the increase in power through an iterative election win. This speaks in favor of
actually awarding the governing party more mandates and not establishing the rule that
all legislative proposals of the governing party are considered adopted as long as only the entire
governing faction agrees --- regardless of its size.

\subsection{Mathematical Questions on Iterative Proportional Representation}

\begin{enumerate}
  \item Does the procedure always terminate?
  \item Must the cut always be made at $\tfrac{2}{3}$, or would another
    ratio also be possible?
\end{enumerate}

Let

\begin{itemize}[label={}]
  \item $A$ -- the total number of votes cast in a round
  \item $A_i$ -- the number of votes cast for the $i$-th party in this round,
    where the parties should be ordered in descending order by their votes. I.e., if
    party $X$ was the $i$-th party in round $j$, it does not
    necessarily have to be so again in round $j+1$.
\end{itemize}

For the relative vote shares $a_i = \frac{A_i}{A}$, the following
conditions apply in each round, where $r$ ($0 < r < 1$) is the relative cutoff vote number and $s$ is the index of the
last party entering the next round:

\begin{itemize}[label={}]
  \item[(I)]   $\displaystyle\sum_{i=1}^{n} a_i = 1$
    \quad ($n$ -- number of participating parties in this round)
  \item[(II)]  $a_k \geq a_l$ \quad for $k < l$
  \item[(III)] $\displaystyle\sum_{i=1}^{s} a_i \geq r$
  \item[(IV)]  $\displaystyle\sum_{i=1}^{s-1} a_i < r$
  \item[\phantom{(IV)}] $\forall i: a_i \neq 0$ \quad
    (Parties without votes are removed from the list)
\end{itemize}

For the procedure to terminate, it must be shown that the number of parties running for election
is strictly reduced in each step. We prove a stronger statement: With $n \geq 3$
parties, at least $\lfloor n(1-r) \rfloor$ parties are eliminated in each step, where
$\lfloor x \rfloor$ denotes the largest integer $\leq x$ (floor function).

\begin{lemma}
For the $k$ weakest parties:
\begin{equation}
  \sum_{i=n-k+1}^{n} a_i \leq \frac{k}{n}
\end{equation}
\end{lemma}

\begin{proof}
  By contradiction. Assumption:
\begin{equation}
  \sum_{i=n-k+1}^{n} a_i > \frac{k}{n}
\end{equation}
Since the total sum is~1, it follows:
\begin{equation}
  \sum_{i=1}^{n-k} a_i = 1 - \sum_{i=n-k+1}^{n} a_i < 1 - \frac{k}{n} = \frac{n-k}{n}
\end{equation}
If all $a_i \geq \frac{1}{n}$ for $i \leq n-k$, their sum would be $\geq \frac{n-k}{n}$,
in contradiction to~(3). Therefore, at least one must be smaller than $\frac{1}{n}$. Since the
$a_i$ are ordered in descending order, it then applies to all following:
\begin{equation}
  a_i < \frac{1}{n} \quad \text{for all } i \geq i^*
\end{equation}
In particular for all $k$ last elements: $\displaystyle\sum_{i=n-k+1}^{n} a_i < \frac{k}{n}$
--- contradiction.
\end{proof}

\begin{maintheorem}
At least $\lfloor n(1-r) \rfloor$ parties are eliminated.
\end{maintheorem}

\begin{proof}
From the lemma, it follows for the remaining $n-k$ parties:
\begin{equation}
  \sum_{i=1}^{n-k} a_i = 1 - \sum_{i=n-k+1}^{n} a_i \geq 1 - \frac{k}{n}
\end{equation}
Condition~(III) is thus satisfied for $s = n-k$ if:
\begin{equation}
  1 - \frac{k}{n} \geq r \quad \Longleftrightarrow \quad k \leq n(1-r)
\end{equation}
Since $k$ must be an integer, at least $\lfloor n(1-r) \rfloor$ parties are eliminated.
\end{proof}

\runinhead{Conclusions for Specific Cutoff Vote Numbers}

\begin{itemize}
  \item \textbf{$r = \frac{1}{2}$:} At least $\lfloor \frac{n}{2} \rfloor$
    parties are eliminated. For $n = 2$, at least $\lfloor 1 \rfloor = 1$ party is eliminated; a
    special rule for $n = 2$ is not necessary.

  \item \textbf{$r = \frac{2}{3}$:} At least $\lfloor \frac{n}{3} \rfloor$
    parties are eliminated. For $n = 3$, at least $\lfloor 1 \rfloor = 1$ party is eliminated.
    For $n = 2$, $\lfloor \frac{2}{3} \rfloor = 0$, i.e., no party is
    guaranteed to be eliminated~-- a special rule for $n = 2$ is therefore necessary.

  \item \textbf{$r = \frac{3}{4}$:} At least $\lfloor \frac{n}{4} \rfloor$
    parties are eliminated. For $n = 2$ and $n = 3$, $\lfloor \frac{n}{4} \rfloor = 0$; a
    special rule is necessary for both cases.
\end{itemize}

\textbf{In general}: The special rule is needed only for $n = 2$~-- and
not also for $n = 3$~-- if and only if $\lfloor 3(1-r) \rfloor \geq 1$, i.e., $r \leq \frac{2}{3}$.
The largest possible value that still satisfies this condition is thus $r = \frac{2}{3}$.

If we were to cut only at three-quarters, for example, we would need a special rule
already for three parties in the round and not just for two.

This can also be easily visualized:

If three parties are in play that all receive exactly the same number of votes, then we cannot
apply the procedure because prerequisite~IV is not satisfied. However, if the parties also
show only minimal vote differences (which can be assumed in elections with many voters and few
candidates), one party will receive less than one-third of the votes and
the other two together will thus exceed two-thirds, meaning only the
two will remain in the next round. But if we were to make the cut at three-quarters, with three
parties, minimal result differences would by no means let a party fall below one-quarter.
The results would have to differ significantly, which can happen but doesn't have to.
If, on the other hand, four parties were in play, minimal result differences would also
cause a party to be eliminated even with a three-quarter cut.

\subsection{Constitutional Amendments}

For good reason, constitutional amendments may only be made with a particularly large majority, the
\enquote{constitutional majority}, which in Germany amounts to two-thirds. Constitutional amendments should only be possible with broad social consensus ---
even across party lines. The two-thirds rule cannot be applied to parliaments with
an awarded absolute majority, as this would feign a broader social
consensus than actually exists. The governing party has received more power in the sense of the
lesser evil than it is actually entitled to, so that at least someone can
decide something, even if ultimately perhaps only the second-best solution is implemented.
At least something is implemented that is coherent.
This principle of the lesser evil should not apply to fundamental questions, however.
Here, correct decisions must be made, as the consequences of wrong decisions would be too serious.
Therefore, I propose that a constitutional chamber be formed that has the composition
as determined by the voter in the first round. All constitutional amendments must be
confirmed with a two-thirds majority in this chamber. If larger cross-party
agreements are necessary for this, then that's how it is and entirely as intended.
The constitutional chamber can also be used beyond constitutional amendments to make other
constitutionally relevant decisions, such as the election of
constitutional judges.
For all constitutionally relevant decisions, it is assumed that the constitution
is actually already good as it is. Therefore, it is deliberately made difficult to change it,
so that it is not accidentally made worse. It was once created through broad
social consensus\footnote{The consensus that led to the formation of the constitution
of the Federal Republic of Germany may not have come about voluntarily, but
through a lost war. However, it has subsequently been accepted in principle by millions of West Germans,
so that one can indeed speak of a society-wide consensus.} and only at least an equally broad social consensus should be able to
change it again. If it cannot be achieved, then it must remain as it is.
Should an equally large consensus not be achieved for decades, then that's how it is and
we continue to live within the framework set by the (intellectual) fathers. The principle only reaches
limits with constitutional judges who, unlike written constitutions,
age over time and must be replaced. Therefore, there must also be a procedure for how new
constitutional judges can be determined who follow in the footsteps of their predecessors, without
the required two-thirds majority necessarily being achieved. The most obvious solution is
actually to let the sitting constitutional judges choose their own successors. They
should not only be able to propose them, but determine them themselves, unless politics intervenes
with a two-thirds majority. On their own, they will appoint successors who will
decide in their spirit.
